\documentclass[11pt,a4paper]{article}
\usepackage{xcolor}
\usepackage[margin=2.5cm]{geometry}
\usepackage{amsmath,amssymb,amsthm,mathtools,bm}
\usepackage{booktabs,graphicx,hyperref,enumitem}
\hypersetup{colorlinks=true,linkcolor=blue!50!black,citecolor=blue!50!black,urlcolor=blue!50!black}
\graphicspath{{../figures/}}
\newtheorem{theorem}{Theorem}[section]
\newtheorem{proposition}[theorem]{Proposition}
\newtheorem{lemma}[theorem]{Lemma}
\theoremstyle{definition}\newtheorem{definition}[theorem]{Definition}
\theoremstyle{remark}\newtheorem{remark}[theorem]{Remark}\newtheorem*{observation}{Numerical observation}
\newcommand{\E}{\mathbb{E}}
\newcommand{\R}{\mathbb{R}}
\DeclareMathOperator{\clip}{clip}

\title{A first-to-default basket credit-linked note:\\ product, pricing and hedging}
\author{Zak Claridge}
\date{October 2026}
\begin{document}
\maketitle
\begin{abstract}
We describe a 5-year funded first-to-default (FtD) credit-linked note (CLN) on five
investment-grade names: how credit default swaps work, why such notes exist, how the FtD spread
depends on default correlation, and how the note is priced (piecewise-constant hazard rates
bootstrapped ISDA-style, one-factor Gaussian copula, semi-analytic and Monte Carlo) and hedged
(single-name CS01, jump-to-default, correlation). All market data is synthetic and illustrative.
\end{abstract}
\tableofcontents

\section*{Notation and conventions}
\begin{itemize}[nosep]
\item Time $t$ in years; premium dates quarterly, $\Delta_i=t_i-t_{i-1}=0.25$ (real CDS use ACT/360 and
IMM dates; we use the same simplified convention in bootstrapping and repricing).
\item $P(t)$: OIS discount factor; $Q_i(t)$: survival probability of name $i$; $\lambda_i$: hazard rate;
$R_i$: recovery rate (40\%); $\tau_i$: default time; $\tau_1=\min_i\tau_i$; $S_1(t)=P(\tau_1>t)$.
\item Protection \emph{buyer} pays the premium and is long credit protection (gains when spreads widen).
The CLN \emph{investor} is short protection. Spreads in basis points (bp) per year.
\end{itemize}

\section{Credit default swaps}
A CDS on a reference entity transfers its credit risk: the protection buyer pays a premium (a spread
or a standard coupon) quarterly until maturity or default; on a \emph{credit event} the seller pays
$1-R$ per unit notional.
\begin{itemize}[nosep]
\item \textbf{Credit events} (ISDA definitions): bankruptcy, failure to pay and, for some contracts
and regions, restructuring (and obligation acceleration/repudiation for some sovereigns). A
Determinations Committee rules whether an event occurred.
\item \textbf{Settlement} is by auction: dealers bid for deliverable bonds and the auction's final
price fixes $R$ for all contracts.
\item \textbf{Standard coupons and upfront.} Since 2009 contracts trade with a fixed coupon (100bp or
500bp for most corporates) and an upfront payment that makes the trade fair:
upfront $=\text{protection}-C\times\mathrm{RPV01}$. A name quoting 125bp with a 100bp coupon has a positive
upfront paid by the buyer.
\item \textbf{Accrued premium.} At default the buyer pays the premium accrued since the last coupon date.
\end{itemize}
With survival $Q$ and recovery $R$, the premium leg per unit spread (RPV01) and the protection leg are
\begin{align*}
\mathrm{RPV01}&=\sum_i\Delta_iP(t_i)Q(t_i)+\sum_i\int_{t_{i-1}}^{t_i}(s-t_{i-1})P(s)\,(-\mathrm dQ(s)),\\
\text{protection}&=(1-R)\int_0^TP(s)\,(-\mathrm dQ(s)),\qquad\text{par spread}=\frac{\text{protection}}{\mathrm{RPV01}}.
\end{align*}
\paragraph{Bootstrap.} Hazards are piecewise constant between the quoted tenors (1, 3, 5, 7, 10 years)
and solved tenor by tenor so that each par spread is reproduced; legs are integrated on a fine grid
with flat forwards between knots, as in the ISDA standard model. Every input is repriced to better
than $10^{-4}$bp. For a flat hazard and flat rate the protection leg is
$(1-R)\frac{\lambda}{\lambda+r}(1-e^{-(\lambda+r)T})$ and $s\approx(1-R)\lambda$ (the ``credit triangle''); both are tests.
Not every set of quotes is attainable with non-negative hazards (e.g.\ 10bp out to 7 years and 800bp at
10 years); the bootstrap rejects such inputs.

\section{The credit-linked note}
\subsection{Structure}
A bankruptcy-remote special purpose vehicle (SPV) issues the note to the investor for $N$ and invests the
cash in high-quality collateral. It enters an FtD swap with a dealer, selling protection on the basket.
The investor receives quarterly coupons of floating (OIS) $+\,s$ until maturity or the first credit event.
At the first credit event the note terminates: the collateral is sold, the dealer is paid $N(1-R_i)$
under the swap, and the investor receives $NR_i$. Without a default the investor receives $N$ at maturity.
Economically, the investor is long the collateral and short first-to-default protection.

\subsection{Why it exists}
\begin{itemize}[nosep]
\item \emph{Investor}: credit exposure in a funded, securities format, without ISDA documentation,
collateral agreements or derivatives permissions, and a coupon well above any single name: about 372bp
over OIS here against single-name 5y spreads of 60--125bp.
\item \emph{Dealer/bank}: offloads concentration or correlation risk (it is long FtD protection) to an
investor, funded, so without counterparty exposure to that investor.
\end{itemize}

\subsection{Why the FtD spread lies between the widest spread and the sum}
\begin{proposition}\label{prop:bounds}
For every $t$, $\max_i P(\tau_i\le t)\le P(\tau_1\le t)\le\sum_iP(\tau_i\le t)$. In the one-factor Gaussian
copula with correlation $\rho\ge0$, $S_1(t)=P(\tau_1>t)$ is non-decreasing in $\rho$; at $\rho=0$ it is
$\prod_iQ_i(t)$, and at $\rho=1$ (comonotone defaults) it is $\min_iQ_i(t)$.
\end{proposition}
\begin{proof}
$\{\tau_i\le t\}\subseteq\{\tau_1\le t\}=\bigcup_j\{\tau_j\le t\}$ gives both bounds. In the copula,
$\{\tau_1>t\}=\bigcap_i\{X_i>c_i(t)\}$ with $c_i(t)=\Phi^{-1}(1-Q_i(t))$ and $(X_i)$ standard Gaussian with
pairwise correlation $\rho$. By Slepian's inequality \cite{slepian1962}, for Gaussian vectors with fixed
unit variances the orthant probability $P(X_i>c_i\ \forall i)$ is non-decreasing in every pairwise
correlation. The extreme cases are independence and $X_1=\dots=X_n$.
\end{proof}
Since spreads are roughly $(1-R)\times$ hazard, the FtD spread moves from about the sum of the spreads
(independence: one of five things going wrong) towards the widest single-name spread (comonotone: the
riskiest name defaults first, and the others with it). Correlation \emph{lowers} the FtD spread: the investor
is long correlation.

\section{Pricing}
\subsection{One-factor Gaussian copula}
Following Li \cite{li2000}, $\tau_i=Q_i^{-1}(\Phi(X_i))$ with $X_i=\sqrt\rho\,M+\sqrt{1-\rho}\,Z_i$. Conditional on $M$
the defaults are independent with $p_i(t|M)=\Phi\big((\Phi^{-1}(1-Q_i(t))-\sqrt\rho M)/\sqrt{1-\rho}\big)$, so
\[
S_1(t)=\int\phi(m)\prod_i\big(1-p_i(t|m)\big)\mathrm dm,
\]
computed by Gauss--Hermite quadrature (Hull--White style semi-analytic pricing \cite{hw2004}; converged to
$10^{-8}$ with 32 nodes). For $\rho>0.9$ the conditional probabilities become steep in $m$ and fixed nodes fail
(at $\rho=0.99$, 64 Gauss--Hermite nodes gave 140bp against a Monte Carlo 128bp); we switch to a fine
trapezoid rule there. With name-specific recoveries the protection leg uses
\[P(\tau_1\in\mathrm dt,\ \text{first}=i\mid M)=\mathrm dp_i(t|M)\prod_{j\ne i}\big(1-p_j(t|M)\big).\]
The FtD swap's legs are those of a CDS with survival $S_1$.

\paragraph{Monte Carlo.} We simulate $\tau_i=\Lambda_i^{-1}(-\log U_i)$, the inverse of the cumulative
hazard, with antithetic variates, and evaluate the legs exactly on each path. The semi-analytic and
Monte Carlo spreads agree within one to two standard errors at every correlation (e.g.\ at $\rho=0.6$,
280.83bp semi-analytic against $281.19\pm0.37$bp from 2 million paths).

\subsection{CLN valuation}
With deterministic rates independent of defaults and coupons stopping at default (no accrued coupon),
\[
\mathrm{PV}=\sum_i\big[P(t_{i-1})-P(t_i)\big]S_1(t_i)+s\sum_i\Delta_iP(t_i)S_1(t_i)
+\int_0^TP(t)\sum_iR_i\,\mathrm dF_{1,i}(t)+P(T)S_1(T),
\]
where the first term values the floating coupons ($P(t_{i-1})/P(t_i)-1$ is the forward rate times
the accrual) and $F_{1,i}$ is the probability that name $i$ is the first to default. A default-free note
is a par floater (PV $=1$ at $s=0$; tested). The fair spread $s^*$ solves PV $=1$.
\begin{observation}
At $\rho=0.3$ the FtD par spread is 366.6bp and the CLN fair spread over OIS is 372.1bp. The difference
arises because the CLN investor receives no accrued coupon at default whereas the FtD protection seller does.
\end{observation}

\subsection{Results}
Five names with 5y spreads 60--125bp (sum 455bp, widest 125bp), recovery 40\%, OIS zero rates 3.6--4.0\%.
The FtD spread falls from 443bp at $\rho=0$ to 130bp at $\rho=0.99$ (Figure in the README); a Student-$t$
copula with 4 degrees of freedom gives lower FtD spreads at every $\rho$ (e.g.\ 297bp vs 339bp at $\rho=0.4$)
because tail dependence makes defaults come together, which helps a first-to-default seller. Even at
$\rho=0$ the $t$ copula has dependence (383.5bp vs 443.4bp), through its common $\chi^2$ mixing variable.

\section{Hedging}
\subsection{Sensitivities}
From the investor's side, at $\rho=0.3$ and the fair spread (per unit notional):
\begin{center}\small
\begin{tabular}{lrrrrr}\toprule
& A & B & C & D & E\\\midrule
5y spread (bp) & 80 & 100 & 60 & 125 & 90\\
CS01 (bp of notional per bp) & $-2.97$ & $-3.08$ & $-2.84$ & $-3.20$ & $-3.03$\\
CS01-neutral 5y CDS hedge (dealer sells) & 0.671 & 0.702 & 0.637 & 0.734 & 0.687\\
Dealer JTD, unhedged & 0.600 & 0.600 & 0.600 & 0.600 & 0.600\\
Dealer JTD, CS01-hedged & 0.197 & 0.179 & 0.218 & 0.159 & 0.188\\\bottomrule
\end{tabular}
\end{center}
CS01 bumps one name's par spread curve by 1bp and re-bootstraps; bucketed CS01 puts almost all the risk
at the 5y node and exactly zero at 7y and 10y. Correl01 is $+11.1$bp of notional per $+1\%$ correlation;
lowering the recovery assumption by 1\% with spreads fixed costs the investor 1.9bp.

\subsection{The dealer's hedge and what it cannot hedge}
The dealer is long FtD protection; it sells 5y protection on each name with notional
$h_i=\mathrm{CS01}^{\rm FtD}_i/\mathrm{CS01}^{\rm CDS}_i$, in total 3.4 times the note notional.
\begin{itemize}[nosep]
\item \textbf{CS01-neutral is not JTD-neutral.} If name $i$ defaults now the dealer gains $1-R=0.6$ on the
FtD but loses only $h_i(1-R)\approx0.4$ on the hedge: it remains long about 0.16--0.22 of jump-to-default
per name. Matching JTD instead ($h_i=1$) would leave large CS01 exposure.
\item \textbf{Correlation is unhedgeable with single-name CDS}: their values do not depend on $\rho$
at all. Correlation risk can be hedged with instruments whose value depends on it, such as index tranches,
where market correlation is quoted as \emph{base correlation}. Mapping a bespoke five-name basket to
index base correlations is itself a model choice.
\item \textbf{Gamma and contagion}: hedge ratios change as spreads move and must be rebalanced; after a
first default the surviving names' spreads typically jump up, which hurts a dealer who is short protection
on them.
\end{itemize}

\subsection{Hedge backtest}
Over one year, hazard curves move by correlated lognormal multipliers (volatility 50\%, correlation
0.5) and a name defaults when its realised cumulative hazard crosses a threshold drawn from the
Gaussian copula ($\rho=0.3$), so spreads widen before defaults. The dealer marks to the model at each
date and rebalances monthly. Over 3{,}000 paths (101 with a default):
\begin{center}\small
\begin{tabular}{lrr}\toprule
One-year P\&L (bp of notional) & unhedged & CS01-hedged\\\midrule
mean (all paths) & $-5$ & $+29$\\
standard deviation (paths without a default) & 485 & 18\\
mean carry & $-365$ & $-48$\\
mean P\&L given a default & $+5{,}871$ & $+1{,}674$\\
5th / 95th percentile & $-881$ / $+1{,}033$ & $-61$ / $-2$\\\bottomrule
\end{tabular}
\end{center}
Hedging removes about 96\% of the spread-driven dispersion. What is left is the default exposure: the
hedged book bleeds a little carry and time decay in normal years and makes about 17\% of notional when a
name defaults. The hedged mean of $+29$bp (standard error about 6bp) is a \emph{model mismatch}, not
alpha. The simulated market has no contagion, so the dealer unwinds its short protection on the survivors
at unchanged spreads after a default. The copula prices those survivors as riskier after a default. The
unhedged position, which has no such hedges, averages zero within error, as fair pricing requires.

\section{Risks}
\begin{itemize}[nosep]
\item \textbf{To the investor:} the first default among five names costs $1-R\approx60\%$ of principal;
correlation (lower correlation hurts the investor); recovery uncertainty (auction outcomes vary widely);
collateral and SPV risk (ignored in the base case); liquidity (bespoke notes are hard to sell); model risk
(the Gaussian copula has no tail dependence).
\item \textbf{To the dealer:} jump-to-default and correlation exposure left after CS01 hedging; gamma and
rebalancing costs; contagion; documentation basis between the note and the hedges.
\end{itemize}

\section{Limitations of this implementation}
Simplified accrual conventions (no ACT/360, IMM dates or business days); deterministic rates independent
of defaults; constant recovery; static copula correlation; no issuer or collateral default; synthetic
market data; the backtest's spread dynamics are not derived from the pricing model, and there is no contagion.

\appendix
\section{Illustrative term sheet (not an offer; synthetic)}
\begin{center}\small
\begin{tabular}{p{0.32\textwidth}p{0.6\textwidth}}\toprule
Issuer & Illustrative SPV Ltd (bankruptcy-remote), collateralised by AAA-rated government bonds\\
Notional & EUR 10,000,000; issue price 100\%\\
Maturity & 5 years\\
Reference basket & Names A--E (five investment-grade corporates; synthetic)\\
Coupon & Compounded overnight index rate + 3.72\% per annum, quarterly, until maturity or the first
credit event (no accrued coupon after a credit event)\\
Credit events & Bankruptcy, failure to pay (ISDA 2014 definitions)\\
Settlement on a credit event & Early redemption at notional $\times$ auction final price of the first
reference entity to suffer a credit event\\
Redemption at maturity & 100\% if no credit event has occurred\\
Main risks & Loss of up to $(1-R)$ of principal on the first default in the basket; correlation, recovery,
collateral, liquidity and model risk\\\bottomrule
\end{tabular}
\end{center}

\begin{thebibliography}{9}
\bibitem{okane2008} D.~O'Kane. \emph{Modelling Single-name and Multi-name Credit Derivatives}. Wiley, 2008.
\bibitem{li2000} D.~X.~Li. On default correlation: a copula function approach. \emph{Journal of Fixed
Income}, 9(4):43--54, 2000.
\bibitem{hw2004} J.~Hull and A.~White. Valuation of a CDO and an $n$th to default CDS without Monte
Carlo simulation. \emph{Journal of Derivatives}, 12(2):8--23, 2004.
\bibitem{isda} ISDA CDS Standard Model, documentation and source code (cdsmodel.com).
\bibitem{slepian1962} D.~Slepian. The one-sided barrier problem for Gaussian noise. \emph{Bell System
Technical Journal}, 41(2):463--501, 1962.
\end{thebibliography}
\end{document}
